\documentclass[dvipsnames]{article}

\usepackage[utf8]{inputenc}
\usepackage{fontspec}
\usepackage{polyglossia}
\usepackage{unicode-math}
\usepackage{csquotes}
\setdefaultlanguage{russian}
\setotherlanguage{english}

\setmainfont{CMU Serif}
\setsansfont{CMU Sans Serif}
\setmonofont{CMU Typewriter Text}
% \setmathfont{Latin Modern Math}

\usepackage[letterpaper,top=2cm,bottom=2cm,left=3cm,right=3cm,marginparwidth=1.75cm]{geometry}

\usepackage{amsmath}
\usepackage{graphicx}
\usepackage{amsthm}
\usepackage{array}
\usepackage{fancyhdr}
\usepackage{bussproofs}
\usepackage[colorlinks=true, allcolors=blue]{hyperref}
\usepackage[backend=biber]{biblatex}
\usepackage{xcolor}

\usepackage{minted2}
\usemintedstyle{xcode}
\setminted{escapeinside=||}
\newmintinline[stella]{c}{fontsize=\scriptsize}
\newmintinline[haskell]{haskell}{fontsize=\scriptsize}
\newmintinline[java]{java}{fontsize=\scriptsize}
\newmintinline[python]{python}{fontsize=\scriptsize}
\newmintinline[code]{c}{fontsize=\small}
\let\codeOrig\code
\renewcommand{\code}[2][]{%
  \ifmmode
    \if\relax\detokenize{#1}\relax
      \mathchoice{\mbox{\normalfont\small\ttfamily #2}}{\mbox{\normalfont\small\ttfamily #2}}{\mbox{\normalfont\small\ttfamily #2}}{\mbox{\normalfont\small\ttfamily #2}}%
    \else
      \mathchoice{\mbox{\normalfont\small\ttfamily #2}}{\mbox{\normalfont\small\ttfamily #2}}{\mbox{\normalfont\small\ttfamily #2}}{\mbox{\normalfont\small\ttfamily #2}}%
    \fi
  \else
    \codeOrig[#1]{#2}%
  \fi}

\usepackage{sourcecodepro}
\usepackage{mdframed}

\DeclareMathOperator{\xor}{\mathsf{xor}}
\DeclareMathOperator{\tru}{\mathsf{tru}}
\DeclareMathOperator{\fls}{\mathsf{fls}}

\addbibresource{../references.bib}

\begin{document}

\pagestyle{fancy}
\fancyhf{}
\chead{Задание выполнил \textcolor{blue}{Павел Григоренко}}
\cfoot{Языки программирования. Семантика и системы типов (весна 2026, Университет ИТМО)}

\section*{Тема 3. Теоретическое задание \textcolor{blue}{(выполнил Павел Григоренко)}}%

В следующих заданиях рассматривается простое типизированное $\lambda$-исчисление с единичным типом, \code{let}-связыванием, парами (кортежами), записями и последовательным исполнением (\code{t|$_1$| ; t|$_2$|}). Напоминание: последовательное исполнение задаётся как производная форма
\[
  \code{t$_1$ ; t$_2$} ::= \code{(λ x : Unit. t$_2$) t$_1$}
\]
с правилами вычисления E-Seq, E-SeqNext и типизации T-Seq~\cite[\S 11.3]{TaPL-ru}.

\begin{enumerate}

  \item Проставьте явные аннотации типа для аргументов $\lambda$-абстракций и постройте дерево вывода типа для замкнутого терма (см. описание пар и записей в~\cite[§§11.6--11.8]{TaPL-ru}):
  \[
  \code{let g = λt. \{ a = t.1, b = t.2 \} in (λf. (f \{ succ 0, iszero (f \{ 0, false \}).a \}).b) g}
  \]

  \color{blue}
  \textbf{Решение.}
  \emph{место для решения}
  \color{black}

  \clearpage
  \item
  Представьте расширение языка выражений \emph{ленивым} логическим оператором ИЛИ: (\code{t|$_1$| or t|$_2$|}).
  Правила вычисления:
  \begin{prooftree}
    \AxiomC{$\code{t$_1$} \longrightarrow \code{t$_1'$}$}
    \RightLabel{E-Or}
    \UnaryInfC{$\code{t$_1$ or t$_2$} \longrightarrow \code{t$_1'$ or t$_2$}$}
  \end{prooftree}
  \begin{prooftree}
    \AxiomC{}
    \RightLabel{E-OrFalse}
    \UnaryInfC{$\code{false or t$_2$} \longrightarrow \code{t$_2$}$}
    \DisplayProof\quad
    \AxiomC{}
    \RightLabel{E-OrTrue}
    \UnaryInfC{$\code{true or t$_2$} \longrightarrow \code{true}$}
  \end{prooftree}
  Правило типизации:
  \begin{prooftree}
    \AxiomC{$\Gamma \vdash \code{t$_1$} : \code{Bool}$}
    \AxiomC{$\Gamma \vdash \code{t$_2$} : \code{Bool}$}
    \RightLabel{T-Or}
    \BinaryInfC{$\Gamma \vdash \code{t$_1$ or t$_2$} : \code{Bool}$}
  \end{prooftree}
  Представьте этот оператор как \emph{производную форму} над простым типизированным $\lambda$-исчислением с логическими и арифметическими выражениями (типы \code{Bool}, \code{Nat}, конструкция \code{if}~\cite[§§8.2--8.3]{TaPL-ru}): выпишите функцию раскрытия сокращений $e$ для термов, содержащих \code{or}, и докажите аналог теоремы~11.3.1~\cite{TaPL-ru} — сохранение вычисления и сохранение типизации при раскрытии.

  \color{blue}
  \textbf{Решение.}
  \emph{место для решения}
  \color{black}

  \clearpage
  \item
  Рассмотрите типы функций:
  \begin{enumerate}
  \item \code{(A × B → C) → ((A → C) × (B → C))};
  \item \code{((A → C) × (B → C)) → (A × B → C)}.
  \end{enumerate}
  Выпишите терм, соответствующий каждому функциональному типу выше, предполагая, что $A$, $B$ и $C$ — неинтерпретируемые базовые типы. Можно использовать возможности типизированного $\lambda$-исчисления с логическими и арифметическими выражениями~\cite[§§8.2--8.3]{TaPL-ru}, единичным типом и последовательным исполнением~\cite[§11.3]{TaPL-ru}, \code{let}-связыванием и приписыванием типа~\cite[§§11.4--11.5]{TaPL-ru}, а также парами, кортежами и записями~\cite[§§11.6--11.8]{TaPL-ru}. Если терм построить невозможно, поясните почему.

  \color{blue}
  \textbf{Решение.}\begin{enumerate}
  \item невозможно
  \item \code{λ acxbc : ((A → C) × (B → C)). λ ab : (A × B). acxbc.1 ab.1} \\
  или \\
  \code{λ acxbc : ((A → C) × (B → C)). λ ab : (A × B). acxbc.2 ab.2}
  \end{enumerate}
  \color{black}

\end{enumerate}

\section*{Список литературы}
\printbibliography[heading=none]

\end{document}
